\documentclass{article}
\usepackage[UTF8]{ctex}
\usepackage{amsmath,amssymb,amsfonts}
\usepackage{enumitem}
\usepackage{geometry}
\geometry{a4paper, margin=1in}
% === 错题本排版优化（不参与题目标识） ===
\setlength{\parindent}{0pt}
\setlength{\parskip}{0.35em plus 0.10em minus 0.08em}
\setlength{\emergencystretch}{3em}
\setlist{itemsep=0.25em, topsep=0.40em, parsep=0pt}
\allowdisplaybreaks[2]
% === 错题本排版优化结束 ===
\begin{document}
\section*{广西大学2025年数学分析真题}
\section*{163．广西大学2025年数学分析试卷}

\subsection{题目1}
1． \(\lim _{x \rightarrow 0} \frac{1}{x} \int_{0}^{x} \cos t^{3} \mathrm{~d} t=\) \(\_\_\_\_\)．

\subsection{题目2}
2．\(f(x)=\frac{3 x^{3}+4}{x^{2}-2 x}\) 的渐近线为 \(\_\_\_\_\)．

\subsection{题目3}
3．幂级数 \(\sum_{n=1}^{\infty} \frac{(x-2)^{2 n-1}}{(2 n-1)!}\) 的收敛区域为 \(\_\_\_\_\)．

\subsection{题目4}
4．\(u=x y z\) 在点（5，1，2）沿点（5，1，2）到 \((9,4,14)\) 的方向导数为 \(\_\_\_\_\)．

\subsection{题目5}
5．\(L\) 是球面 \(x^{2}+y^{2}+z^{2}=a^{2}\) 与 \(x=y\) 相交的圆周，则 \(\int_{L} \sqrt{2 y^{2}+z^{2}} \mathrm{~d} s=\) \(\_\_\_\_\)．

\subsection{题目6}
6．求 \(y=x^{2}\) 与 \(y=8-x^{2}\) 所围区域的面积．

\subsection{题目7}
7．\(u=x^{2}+y^{2}+z^{2}\) ，其中 \(z=f(x, y)\) 是由 \(x^{3}+y^{3}+z^{3}=3 x y z\) 确定的隐函数，求 \(u_{x}, u_{x x}\) ．

\subsection{题目8}
8．设 \(f(x)=\left\{\begin{array}{ll}-\frac{\pi}{4}, & -\pi<x<0 ; \\ \frac{\pi}{4}, & 0 \leq x<\pi .\end{array}\right.\) 求 \(f(x)\) 的傅里叶级数展开式．并推出

\[
\frac{\sqrt{3}}{6} \pi=1-\frac{1}{5}+\frac{1}{7}-\frac{1}{11}+\frac{1}{13}-\frac{1}{17}+\cdots .
\]

\subsection{题目9}
9 ．求 \(\left(x_{0}, y_{0}, z_{0}\right)\) 到平面 \(A x+B y+C z+D=0\) 最短距离．

\subsection{题目10}
10．求 \(I=\iiint_{V} z^{2} \mathrm{~d} x \mathrm{~d} y \mathrm{~d} z, V\) 是由锥面 \(z^{2}=\frac{h^{2}}{R^{2}}\left(x^{2}+y^{2}\right)\) 与平面 \(z=h\) 围成的区域．

\subsection{题目11}
11．求曲面积分

\[
I=\iint_{S} y(x-z) \mathrm{d} y \mathrm{~d} z+x^{2} \mathrm{~d} z \mathrm{~d} x+\left(y^{2}+z x\right) \mathrm{d} x \mathrm{~d} y
\]

其中 \(S\) 是曲面 \(z=5-x^{2}-y^{2}\) 上 \(z \geq 1\) 部分的外侧。

\subsection{题目12}
12． \(\lim _{n \rightarrow \infty} a_{n}=a, a>0\) ，证明： \(\lim _{n \rightarrow \infty} \sqrt[n]{a_{n}}=1\) ．

\subsection{题目13}
13．证明 \(f(x)=x^{2}\) 在 \([a, b]\) 上一致连续，在 \((-\infty,+\infty)\) 上不一致连续．

\subsection{题目14}
14．设曲线

\[
\left\{\begin{array}{l}
x=2(\sin t-t \cos t) \\
y=2(\cos t+t \sin t)
\end{array}\right.
\]

证明原点到曲线上任意一点处的法线的距离为常数．

\subsection{题目15}
15．根据有限覆盖定理证明：实数轴上的任一有界无限点集至少有一个聚点．

\subsection{题目16}
16．用

\[
f(x, y)= \begin{cases}\left(x^{2}+y^{2}\right) \sin \frac{1}{\sqrt{x^{2}+y^{2}}}, & x^{2}+y^{2} \neq 0 \\ 0, & x^{2}+y^{2}=0\end{cases}
\]

说明偏导数连续并不是函数可微的必要条件。

\subsection{题目17}
17．用求导，积分计算 \(I(x)=\int_{0}^{+\infty} e^{-t^{2}} \cos 2 x t \mathrm{~d} t\) ，并说明原因．
\end{document}
